\SourcePage{413}{416}
\section{Interaction of atoms at large separations}

Consider two atoms separated by a distance large compared with their dimensions, and determine their interaction energy. In other words, we seek the form of the electronic terms at large internuclear distances.

We apply perturbation theory, regarding the two isolated atoms as the unperturbed system and the potential energy of their electrical interaction as the perturbation operator. As is known (see Vol.~II, \S\S~41, 42), the electrical interaction of two systems of charges separated by a large distance $r$ can be expanded in powers of $1/r$. Successive terms describe the interaction of the total charges, dipole moments, quadrupole moments, and so forth of the two systems. For neutral atoms the total charges vanish. The expansion then begins with the dipole--dipole interaction ($\sim1/r^3$), followed by dipole--quadrupole terms ($\sim1/r^4$), quadrupole--quadrupole (and dipole--octupole) terms ($\sim1/r^5$), and so on.

Suppose first that both atoms are in $S$ states. It is readily seen that their interaction has no effect in first-order perturbation theory. Indeed, to first order the interaction energy is the diagonal matrix element of the perturbation operator evaluated with the unperturbed wave functions of the system (which are themselves products of the wave functions of the two atoms).\footnote{Exchange effects that decrease exponentially with distance (compare Problem~1 of \S~62 and the problem in \S~81) are discarded here.} But in $S$ states the diagonal matrix elements, that is, the mean dipole, quadrupole, and higher moments of the atoms, vanish directly as a consequence of the spherical symmetry of the atomic charge distributions. Hence every term in the expansion of the perturbation operator in powers of $1/r$ gives zero in first-order perturbation theory.\footnote{This does not, of course, mean that the mean interaction energy of the atoms is exactly zero. It decreases exponentially with distance, faster than every finite power of $1/r$, which is why every term in the expansion vanishes. The multipole expansion of the interaction operator assumes that the charges of the two atoms are separated by the large distance $r$. Yet the quantum-mechanical electron-density distribution has finite, though exponentially small, values even at large distances.}

\SourcePage{414}{417}
In second order it is enough to retain the dipole interaction in the perturbation operator, since it decreases most slowly as $r$ increases, namely
\begin{equation}
V=\frac{\mathbf d_1\mathbf d_2-3(\mathbf d_1\mathbf n)(\mathbf d_2\mathbf n)}{r^3}
\tag{89.1}
\end{equation}
($\mathbf n$ is the unit vector directed from atom 1 to atom 2). Since the nondiagonal matrix elements of the dipole moment are in general nonzero, second-order perturbation theory gives a nonzero result which, being quadratic in $V$, is proportional to $1/r^6$. The second-order correction to the lowest eigenvalue is always negative (see \S~38). Thus for atoms in their normal states the interaction energy has the form
\begin{equation}
U(r)=-\frac{\mathop{\rm const}\nolimits}{r^6},
\tag{89.2}
\end{equation}
where $\mathop{\rm const}\nolimits$ is positive\footnote{For example, the values of this constant (in atomic units) for pairs of atoms are: hydrogen, 6.5; helium, 1.5; argon, 68; krypton, 130.} (F.~London, 1928).

Consequently two atoms in normal $S$ states, at a large separation, attract with a force ($-dU/dr$) inversely proportional to the seventh power of the distance. Attractive forces between atoms at large distances are usually called \emph{van der Waals forces}. These forces also produce wells in the potential-energy curves of electronic terms for atoms that do not form a stable molecule. Such wells are very shallow, however (their depths are only tenths or even hundredths of an electron volt), and occur at distances several times larger than the interatomic distances in stable molecules.

If only one atom is in an $S$ state, the same result (89.2) is obtained for their interaction energy, since the vanishing of the dipole and higher moments of just one atom is enough to make the first-order result zero. The constant in the numerator of (89.2) then depends not only on the states of both atoms but also on their relative orientation, that is, on the magnitude $\Omega$ of the projection of angular momentum on the line joining them.

If both atoms have nonzero orbital and total angular momenta, the situation changes. The mean dipole moment is zero in every atomic state (see \S~75). The mean quadrupole moments (in states with $L\ne0$, $J\ne0,1/2$), however, are nonzero.

\SourcePage{415}{418}
The quadrupole--quadrupole term in the perturbation operator therefore contributes already in first order, and the interaction energy decreases as the fifth rather than the sixth power of the distance:
\begin{equation}
U(r)=\frac{\mathop{\rm const}\nolimits}{r^5}.
\tag{89.3}
\end{equation}
The constant may have either sign, so either attraction or repulsion may occur. As before, it depends not only on the atomic states but also on the state of the combined system formed by the two atoms.

A special case is the interaction between two identical atoms in different states. The unperturbed system (the two isolated atoms) then has an additional degeneracy associated with interchanging the states between the atoms. Accordingly, the first-order correction is determined by a secular equation containing both diagonal and nondiagonal matrix elements of the perturbation. If the two atomic states have opposite parities and angular momenta $L$ differing by $\pm1$ or 0, but not both zero (and likewise for $J$), the nondiagonal dipole-moment matrix elements for transitions between them are in general nonzero. A first-order effect then arises already from the dipole term in the perturbation operator. The atomic interaction energy is consequently proportional to $1/r^3$:
\begin{equation}
U(r)=\frac{\mathop{\rm const}\nolimits}{r^3};
\tag{89.4}
\end{equation}
the constant may have either sign.

Usually, however, the interaction averaged over all possible orientations of the atomic angular momenta is of interest (as, for example, for atoms interacting in a gas). This averaging makes the mean values of all multipole moments vanish. With them vanish all first-order effects linear in these moments in the atomic interaction. Hence in every case the averaged long-range forces between atoms obey (89.2).\footnote{This law, obtained from nonrelativistic theory, applies only while retardation of the electromagnetic interactions is unimportant. For this purpose the separation $r$ must be small compared with $c/\omega_{0n}$, where $\omega_{0n}$ are the transition frequencies between the ground and excited atomic states. For the interaction including retardation, see Vol.~IV, \S~85.}

\SourcePage{416}{419}
We conclude with the related question of the interaction between a neutral atom and an ion.

In first-order perturbation theory this interaction is the mean value of the operator (76.8), the energy of a quadrupole in the Coulomb field of the ion. Since the ion's potential $\varphi\sim1/r$, the atom--ion interaction energy is proportional to $1/r^3$. This effect exists only if the atom has a nonzero mean quadrupole moment. Even then, it disappears on averaging over all directions of the atomic angular momentum $\mathbf J$.

The next term in powers of $1/r$, which is always nonzero, is the second-order interaction from the dipole operator (76.1). Since the ion's field strength is $\sim1/r^2$, this interaction energy is proportional to $1/r^4$. In terms of the polarizability $\alpha$ of an atom in an $S$ state, it is
\begin{equation}
U=-\alpha e^2/2r^4.
\tag{89.5}
\end{equation}
If the atom is in its normal state, this energy (like every correction to the ground-state energy) is negative, so an attractive force acts between atom and ion.\footnote{The same attraction occurs at large distances between an atom and an electron. It is the reason atoms can form negative ions by attaching an electron (with binding energies from fractions to several electron volts). Not every atom has this property, however. In a field decreasing at large distances as $1/r^4$ (or $1/r^3$), the number of levels corresponding to bound electron states is finite in every case and may, in particular cases, be zero.}

\ProblemHeading For two identical atoms in $S$ states, derive a formula expressing the van der Waals forces in terms of matrix elements of their dipole moments.

\SolutionHeading The result follows from the general perturbation-theory formula (38.10), applied to the operator (89.1). Because atoms in $S$ states are isotropic, it is clear in advance that, on summing over all intermediate states, the squared matrix elements of the three components of each vector $\mathbf d_1$ and $\mathbf d_2$ make equal contributions, whereas terms containing products of different components vanish.

The result is
\[
U(r)=-\frac6{r^6}\sum_{n,n'}\frac{|\langle n|d_z|0\rangle|^2|\langle n'|d_z|0\rangle|^2}{E_n+E_{n'}-E_0},
\]

\SourcePage{417}{420}
where $E_0,E_n$ are the unperturbed energies of the ground and excited atomic states. Since, by assumption, $L=0$ in the ground state, the matrix elements $(d_z)_{0n}$ are nonzero only for transitions to $P$ states ($L=1$). With formulas (29.7), we rewrite $U(r)$ finally as
\[
U(r)=-\frac2{3r^6}\sum_{n,n'}\frac{\langle n1\|d\|00\rangle^2\langle n'1\|d\|00\rangle^2}{E_{n1}+E_{n'1}-2E_{00}},
\]
where in the indices $nL$ of the energy levels and reduced matrix elements, the second symbol gives $L$, while the first denotes the set of all remaining quantum numbers specifying the energy level.
